\hypertarget{advanced-move-semantics-value-categories}{%
\chapter{ADVANCED MOVE SEMANTICS \& VALUE
CATEGORIES}\label{advanced-move-semantics-value-categories}}

``Move Semantics'' is often misunderstood. It's not magic; it's type
casting.

\hypertarget{the-c17-value-category-taxonomy}{%
\subsection{4.5.1 The C++17 Value Category
Taxonomy}\label{the-c17-value-category-taxonomy}}

Everything in C++ is an Expression, and every expression has a
\textbf{Type} and a \textbf{Value Category}.

\begin{lstlisting}
        Expression
       /          \
   glvalue      rvalue
   /    \      /     \
lvalue   xvalue   prvalue
\end{lstlisting}

\begin{enumerate}
\def\labelenumi{\arabic{enumi}.}
\tightlist
\item
  \textbf{lvalue (Identity + Movable? No)}: Has a name, persists beyond
  expression.

  \begin{itemize}
  \tightlist
  \item
    \passthrough{\lstinline!int x; x!} is an lvalue.
  \item
    \passthrough{\lstinline!std::string s; s!} is an lvalue.
  \end{itemize}
\item
  \textbf{prvalue (Pure Rvalue)}: No name, temporary, initializes an
  object.

  \begin{itemize}
  \tightlist
  \item
    \passthrough{\lstinline!10!}, \passthrough{\lstinline!true!},
    \passthrough{\lstinline!nullptr!}.
  \item
    \passthrough{\lstinline!std::string("hello")!} (constructor call).
  \end{itemize}
\item
  \textbf{xvalue (eXpiring Value)}: Has identity, but can be moved from.

  \begin{itemize}
  \tightlist
  \item
    Result of \passthrough{\lstinline!std::move(x)!}.
  \item
    Rvalue reference cast
    \passthrough{\lstinline!static\_cast<T\&\&>(x)!}.
  \end{itemize}
\end{enumerate}

\textbf{glvalue} = lvalue + xvalue (Has Identity) \textbf{rvalue} =
prvalue + xvalue (Can be moved from)

\hypertarget{stdmove-and-stdforward-internals}{%
\subsection{4.5.2 std::move and std::forward
Internals}\label{stdmove-and-stdforward-internals}}

\textbf{\passthrough{\lstinline!std::move!}}: Does NOT move. It
unconditionally casts to rvalue reference.

\begin{lstlisting}[language={C++}]
template<typename T>
typename remove_reference<T>::type&& move(T&& t) noexcept {
    return static_cast<typename remove_reference<T>::type&&>(t);
}
\end{lstlisting}

\textbf{\passthrough{\lstinline!std::forward!}}: Conditionally casts to
rvalue reference \emph{only if} the argument was initialized with an
rvalue. Used for \textbf{Perfect Forwarding}.

\begin{lstlisting}[language={C++}]
template<typename T>
T&& forward(typename remove_reference<T>::type& t) noexcept {
    return static_cast<T&&>(t);
}
\end{lstlisting}

\hypertarget{reference-collapsing-rules}{%
\subsection{4.5.3 Reference Collapsing
Rules}\label{reference-collapsing-rules}}

When templates meet references, types collapse:

\begin{itemize}
\tightlist
\item
  \passthrough{\lstinline!T\& \&!} -\textgreater{}
  \passthrough{\lstinline!T\&!}
\item
  \passthrough{\lstinline!T\& \&\&!} -\textgreater{}
  \passthrough{\lstinline!T\&!}
\item
  \passthrough{\lstinline!T\&\& \&!} -\textgreater{}
  \passthrough{\lstinline!T\&!}
\item
  \passthrough{\lstinline!T\&\& \&\&!} -\textgreater{}
  \passthrough{\lstinline!T\&\&!} (The only way to get an rvalue
  reference)
\end{itemize}

This is why \passthrough{\lstinline!T\&\&!} in a template is a
\textbf{Universal Reference} (Forwarding Reference). It can become
\passthrough{\lstinline!T\&!} (lvalue) or
\passthrough{\lstinline!T\&\&!} (rvalue).

\begin{center}\rule{0.5\linewidth}{0.5pt}\end{center}
